\section*{Erd\H{o}s Problem 565}

\paragraph{Statement and scope.}
The induced Ramsey number \(R^*(G)\) is the least order of a graph
\(H\) whose every red-blue edge colouring contains a monochromatic
copy of \(G\) induced in \(H\). Must \(R^*(G)\leq2^{O(n)}\)
hold for every \(n\)-vertex graph \(G\)?
The general affirmative theorem is known. This attempt proves the
exponential scale is necessary, gives a full exponential upper bound
for complete targets, and supplies small induced hosts for stars and
edgeless targets. The general construction is unresolved here.

\paragraph{Complete targets reduce exactly to ordinary Ramsey numbers.}
Write \(R(n,n)\) for the ordinary two-colour clique Ramsey number.
The complete host \(K_{R(n,n)}\) shows
\(R^*(K_n)\leq R(n,n)\): any monochromatic clique is automatically
induced as an uncoloured clique.

Conversely, if \(N<R(n,n)\), choose a colouring of \(K_N\) with
no monochromatic \(K_n\). Restrict it to the edge set of any
\(N\)-vertex host \(H\). This restriction still has no monochromatic
\(K_n\), so no such host is induced Ramsey for \(K_n\). Thus
\[
 R^*(K_n)=R(n,n). \tag{1}
\]
This argument checks the quantifier over all possible host graphs.

The elementary recurrence
\(R(a,b)\leq R(a-1,b)+R(a,b-1)\) follows by splitting the
neighbours of one vertex according to edge colour.
With the base cases \(R(1,b)=R(a,1)=1\), induction using Pascal's
identity gives
\[
 R(a,b)\leq\binom{a+b-2}{a-1}.
\]
Consequently
\[
 R^*(K_n)\leq\binom{2n-2}{n-1}\leq4^{n-1}. \tag{2}
\]

\paragraph{A self-contained exponential lower bound.}
Colour each edge of \(K_N\) independently red or blue with equal
probability. A specified \(n\)-set is monochromatic with probability
\(2^{1-\binom n2}\). Therefore the expected number of monochromatic
\(K_n\)'s is
\[
 \binom Nn2^{1-\binom n2}.
\]
For \(N=\lfloor2^{n/2}\rfloor\) this is at most
\(2^{1+n/2}/n!<1\) for every \(n\geq4\).
The last inequality holds at \(n=4\), and its left-to-right ratio
decreases when \(n\) is increased because the factorial gains a
factor \(n+1>\sqrt2\).
Some colouring therefore contains no monochromatic \(K_n\).
By (1),
\[
 R^*(K_n)>\lfloor2^{n/2}\rfloor. \tag{3}
\]
Thus any bound valid for every \(n\)-vertex target must allow
exponential growth. A polynomial general upper bound is impossible.

\paragraph{Two other target families with induced hosts.}
For the star on \(n\geq2\) vertices, use a host star with
\(2n-3\) leaves. At least \(n-1\) edges at its centre have the
same colour. Their endpoints induce exactly the desired star,
because the host has no edges between leaves. Hence
\[
 R^*(K_{1,n-1})\leq2n-2. \tag{4}
\]
For the edgeless \(n\)-vertex graph, the edgeless host of the same
order works, with monochromaticity vacuous; fewer than \(n\)
vertices cannot contain the target. Its induced Ramsey number is
therefore exactly \(n\).

\paragraph{The obstruction to extending these arguments.}
For a general non-complete target, the complete host used in (2)
does not work: it contains no induced copy of any graph with a
missing edge. The host must simultaneously retain the prescribed
nonedges and force all its required edges into one colour.
Equations (1)--(4) do not provide such a construction for arbitrary
targets, and so do not prove the main assertion.

The full theorem of Arag\~ao, Campos, Dahia, Filipe and Marciano
gives \(R_{\mathrm{ind}}(G;r)\leq r^{Crn}\) with an absolute
constant \(C\); setting \(r=2\) settles the question.
Its general embedding argument is not reconstructed or used to
fill this attempt's gap. Primary source:
\url{https://arxiv.org/abs/2509.22629}.
Current statement: \url{https://www.erdosproblems.com/565},
checked 24 September 2026. The elementary results above are standard;
no novelty is claimed.

\paragraph{Completion estimate.}
The necessary exponential scale and three target families are proved;
the upper bound for arbitrary targets remains missing from this attempt.
\textbf{COMPLETION: 30\%}
